\paragrafos{Εφεξής και διαδοχικές γωνίες: οι γωνίες που προστίθενται}{Β.1.7, «Εφεξής και διαδοχικές γωνίες – Άθροισμα γωνιών»}

\begin{anagke}
Στο κέντρο του κήπου θα υπάρχει το συντριβάνι και από αυτό θα ξεκινούν πέντε μονοπάτια προς διαφορετικές κατευθύνσεις. Η Αγγελική σχεδιάζει το σχήμα και μετράει τις γωνίες ανάμεσα στα διαδοχικά μονοπάτια: $60^\circ$, $75^\circ$, $50^\circ$, $90^\circ$. «Η πέμπτη γωνία δεν χωρούσε στο μοιρογνωμόνιο», λέει, «γιατί είναι μη κυρτή;». «Όχι», απαντάει η Πολυξένη, «απλώς πρέπει να βρούμε πόσο μένει μέχρι τον κύκλο. Ολόκληρος ο κύκλος έχει $360^\circ$».

Ο Βασίλης προσθέτει: $60+75+50+90=275$, άρα η πέμπτη είναι $360-275=85$. «Μπορεί να μην βρίσκομαι εκεί, αλλά είμαι σίγουρος!». Ο Στέλιος ρωτάει αν οι γωνίες που έχουν κοινή πλευρά και «χωρίζουν» το μέρος που έχουν κοινή κορυφή έχουν κάποιο όνομα. Η κυρία Αλεξάνδρα απαντάει: «Λέγονται εφεξής, και όταν είναι πολλές στη σειρά, διαδοχικές. Τα μέτρα τους προστίθενται».

Στην παράγραφο αυτή μαθαίνουμε πώς προσθέτουμε και αφαιρούμε γωνίες που έχουν κοινή κορυφή και κοινή πλευρά, και πώς χρησιμοποιούμε το γεγονός ότι γύρω από ένα σημείο υπάρχουν $360^\circ$ και πάνω σε μια ευθεία $180^\circ$.
\end{anagke}

\begin{lexi}{εφεξής, διαδοχικές}{εφεξής, διαδοχικές}
\lexiword{Εφεξής} σημαίνει «η μία μετά την άλλη, στη σειρά» (εφ' εξής). \lexiword{Διαδοχικές} είναι οι γωνίες που «διαδέχονται» η μία την άλλη, έτσι ώστε η κάθε μία να είναι εφεξής με την επόμενη.
\end{lexi}

\begin{metafrasi}
Δύο γωνίες με κοινή κορυφή και κοινή πλευρά, και χωριστά εσωτερικά. & εφεξής γωνίες $x\mathrm{Ο}y$, $y\mathrm{Ο}z$\\
Το άθροισμά τους. & $x\mathrm{Ο}z=x\mathrm{Ο}y+y\mathrm{Ο}z$\\
Όλες οι διαδοχικές γύρω από ένα σημείο. & άθροισμα $360^\circ$\\
Οι διαδοχικές πάνω σε ευθεία. & άθροισμα $180^\circ$\\
\end{metafrasi}

\tmima{Η θεωρία}

\begin{orismos}{εφεξής και διαδοχικές γωνίες}
Δύο γωνίες $x\mathrm{Ο}y$ και $y\mathrm{Ο}z$ λέγονται \textbf{εφεξής} όταν έχουν κοινή κορυφή $\mathrm{Ο}$, κοινή πλευρά $\mathrm{Οy}$ και τα εσωτερικά τους δεν έχουν κοινά σημεία (η μία είναι «πριν» και η άλλη «μετά» την $\mathrm{Οy}$). Πολλές γωνίες λέγονται \textbf{διαδοχικές} όταν καθεμία είναι εφεξής με την επόμενή της.
\end{orismos}

\begin{center}
\begin{tikzpicture}[scale=1.1]
\draw[thick] (3.4,0)--(0,0)--(2.4,2.4); \draw[thick,gblue] (0,0)--(0.6,3.1); \fill (0,0) circle (2pt);
\node[right] at (3.4,0) {$x$}; \node[right] at (2.4,2.4) {$y$}; \node[above] at (0.6,3.1) {$z$}; \node[below left] at (0,0) {$\mathrm{Ο}$};
\draw (0.8,0) arc(0:45:0.8); \draw[gblue] (0.55,0.55) arc(45:79:0.55);
\end{tikzpicture}
\end{center}

\tmimasmall{Άθροισμα γωνιών}

Αν οι $x\mathrm{Ο}y$ και $y\mathrm{Ο}z$ είναι εφεξής, η γωνία $x\mathrm{Ο}z$ λέγεται \textbf{άθροισμά} τους:
$$x\mathrm{Ο}z=x\mathrm{Ο}y+y\mathrm{Ο}z.$$
Και αντίστροφα η διαφορά: $y\mathrm{Ο}z=x\mathrm{Ο}z-x\mathrm{Ο}y$. Το άθροισμα των μέτρων μπορεί να ξεπεράσει τα $180^\circ$ (τότε η γωνία είναι μη κυρτή).

\begin{idiotita}{άθροισμα διαδοχικών γωνιών}
Το άθροισμα των διαδοχικών γωνιών γύρω από ένα σημείο, που καλύπτουν όλο το επίπεδο, είναι $360^\circ$. Το άθροισμα διαδοχικών γωνιών που βρίσκονται από την ίδια μεριά μιας ευθείας (πλευρές αντικείμενες ημιευθείες) είναι $180^\circ$.
\end{idiotita}

\begin{prosoxi}
\begin{itemize}[leftmargin=1.2em,itemsep=0.3ex,topsep=2pt]
\item Δύο γωνίες δεν είναι εφεξής αν δεν έχουν κοινή κορυφή ή αν τα εσωτερικά τους επικαλύπτονται.
\item Το άθροισμα γωνιών αθροίζει τα \emph{μέτρα}: $40^\circ+65^\circ=105^\circ$.
\item Στον κύκλο, όλο το άθροισμα είναι $360^\circ$, όχι $180^\circ$ (αυτό είναι για την ευθεία).
\item Για εφεξής γωνίες το άθροισμα μπορεί να ξεπεράσει τις $180^\circ$ αλλά η γωνία $x\mathrm{Ο}z$ είναι τότε μη κυρτή.
\end{itemize}
\end{prosoxi}

\tmima{Λυμένα παραδείγματα}

\paradeigma{Εφεξής γωνίες}
\tixero $x\mathrm{Ο}y=35^\circ$ και $y\mathrm{Ο}z=50^\circ$ εφεξής.\par
\tithelo Η $x\mathrm{Ο}z$.\par
\possindeo $x\mathrm{Ο}z=35^\circ+50^\circ=85^\circ$.\par
\apantisi $85^\circ$.\par
\elegxos $85^\circ-35^\circ=50^\circ$.

\paradeigma{Με αφαίρεση}
\tixero Οι $\mathrm{ΑΟΒ}$, $\mathrm{ΒΟΓ}$ εφεξής με $\mathrm{ΑΟΒ}=70^\circ$ και $\mathrm{ΑΟΓ}=130^\circ$.\par
\tithelo Η $\mathrm{ΒΟΓ}$.\par
\possindeo $\mathrm{ΒΟΓ}=\mathrm{ΑΟΓ}-\mathrm{ΑΟΒ}=130^\circ-70^\circ=60^\circ$.\par
\apantisi $60^\circ$.\par
\elegxos $70^\circ+60^\circ=130^\circ$.

\paradeigma{Τα πέντε μονοπάτια}
\tixero Τέσσερις διαδοχικές γωνίες γύρω από το συντριβάνι: $60^\circ$, $75^\circ$, $50^\circ$, $90^\circ$.\par
\tithelo Η πέμπτη γωνία.\par
\possindeo Το άθροισμα γύρω από το σημείο είναι $360^\circ$. Τέσσερις γωνίες: $60+75+50+90=275$. Άρα η πέμπτη είναι $360-275=85^\circ$.\par
\apantisi $85^\circ$.\par
\elegxos $275^\circ+85^\circ=360^\circ$.

\paradeigma{Γωνίες πάνω σε ευθεία}
\tixero Διαδοχικές γωνίες $40^\circ$, $65^\circ$ και $\omega$ πάνω από μια ευθεία (η τρίτη κλείνει στην ευθεία).\par
\tithelo Η $\omega$.\par
\possindeo Το άθροισμα είναι $180^\circ$: $40+65+\omega=180$, άρα $\omega=180-105=75^\circ$.\par
\apantisi $75^\circ$.\par
\elegxos $40+65+75=180$.

\paradeigma{Διχοτόμος και άθροισμα}
\tixero Η $\mathrm{Οδ}$ είναι διχοτόμος της $x\mathrm{Ο}y$ και $\delta\mathrm{Ο}y=37^\circ$. Η $y\mathrm{Ο}z=48^\circ$ είναι εφεξής στην $x\mathrm{Ο}y$.\par
\tithelo Η $x\mathrm{Ο}z$.\par
\possindeo $x\mathrm{Ο}y=2\cdot37^\circ=74^\circ$. $x\mathrm{Ο}z=74^\circ+48^\circ=122^\circ$.\par
\apantisi $122^\circ$.\par
\elegxos Η $\mathrm{Οδ}$ πλευρά της $x\mathrm{Ο}z$: $37+37+48=122$.

\tmima{Ασκήσεις}

\zoni{Το κατάλαβα}{άθροισμα και διαφορά γωνιών}

\ex Οι $x\mathrm{Ο}y=35^\circ$ και $y\mathrm{Ο}z=50^\circ$ είναι εφεξής. Βρες την $x\mathrm{Ο}z$.

\ex Οι $\mathrm{ΑΟΒ}$ και $\mathrm{ΒΟΓ}$ είναι εφεξής, $\mathrm{ΑΟΓ}=130^\circ$ και $\mathrm{ΑΟΒ}=70^\circ$. Βρες την $\mathrm{ΒΟΓ}$.

\ex Τέσσερις διαδοχικές γωνίες γύρω από σημείο είναι $60^\circ$, $75^\circ$, $50^\circ$, $90^\circ$. Βρες την πέμπτη.

\ex Οι $40^\circ$, $65^\circ$ και $\omega$ είναι διαδοχικές πάνω από ευθεία. Βρες την $\omega$.

\ex Τρεις διαδοχικές γωνίες γύρω από σημείο είναι ίσες. Πόσες μοίρες είναι η καθεμία;

\ex Οι γωνίες $x\mathrm{Ο}y=24^\circ30'$ και $y\mathrm{Ο}z=37^\circ45'$ είναι εφεξής. Βρες την $x\mathrm{Ο}z$.

\ex Σωστό ή λάθος; α) Δύο γωνίες με κοινή κορυφή είναι πάντα εφεξής. \quad β) Οι εφεξής γωνίες έχουν κοινή πλευρά. \quad γ) Διαδοχικές γωνίες γύρω από σημείο έχουν άθροισμα $360^\circ$. \quad δ) Διαδοχικές γωνίες πάνω από ευθεία έχουν άθροισμα $180^\circ$.

\zoni{Το χρησιμοποιώ}{κήπος, πίτσα, ρολόι}

\ex \textbf{Τα πέντε μονοπάτια.} Στο συντριβάνι καταλήγουν πέντε μονοπάτια και οι γωνίες είναι $60^\circ$, $75^\circ$, $50^\circ$, $90^\circ$ και $\omega$. Βρες την $\omega$.

\ex \textbf{Η πίτσα.} Μια πίτσα κόβεται σε $6$ φέτες: δύο από $75^\circ$ και τρεις από $50^\circ$. Πόσο είναι η έκτη;

\ex \textbf{Ο ήλιος.} Ο Ήλιος σηκώνεται από τον ορίζοντα $12^\circ$ κάθε $48$ λεπτά. Σε πόση ώρα θα φτάσει να ανεβεί $60^\circ$;

\ex \textbf{Το ρολόι.} Στις $4{:}30$ ο ωροδείκτης βρίσκεται στις $135^\circ$ από το $12$ και ο λεπτοδείκτης στο $6$ ($180^\circ$ από το $12$). Πόση γωνία σχηματίζουν;

\ex \textbf{Οι στίβοι.} Τρεις στίβοι ξεκινούν από το ίδιο σημείο σχηματίζοντας ίσες γωνίες. Πόση είναι η γωνία ανάμεσα σε δύο διαδοχικούς;

\zoni{Το σκέφτομαι}{όταν το σύνολο είναι σταθερό}

\ex Γιατί το άθροισμα των διαδοχικών γωνιών γύρω από σημείο είναι $360^\circ$ ανεξάρτητα από τον αριθμό των γωνιών;

\ex Είναι δυνατόν δύο εφεξής γωνίες να είναι και οι δύο αμβλείες; Εξήγησε.

\ex Τρεις διαδοχικές γωνίες γύρω από σημείο είναι $x$, $2x$, $3x$. Βρες το $x$.

\ex Η $\mathrm{Οδ}$ είναι διχοτόμος της $x\mathrm{Ο}y=60^\circ$ και η $y\mathrm{Ο}z=20^\circ$ είναι εφεξής προς αυτήν. Βρες τις $x\mathrm{Ο}z$ και $\delta\mathrm{Ο}z$.

\zoni{Για το σπίτι}{άθροισμα στο χαρτί}

\ex Σχεδίασε τέσσερις διαδοχικές γωνίες γύρω από σημείο και μέτρησέ τις. Πόσο είναι το άθροισμά τους;

\ex Βρες στο σπίτι ένα αντικείμενο με διαδοχικές γωνίες γύρω από σημείο (π.χ. ένα ρολόι ή μια ανεμογεννήτρια).
